\documentclass[a4paper]{article}
\usepackage[pdftex]{graphicx}
\usepackage[utf8]{inputenc}
\usepackage{enumerate}
\usepackage{icomma}
\usepackage{amssymb}
\usepackage{tikz}
\usepackage{href-ul}
\hypersetup{
	colorlinks=true,
	linkcolor=blue,
	urlcolor=blue}
\usepackage{geometry}
\geometry{a4paper, top=15mm, left=15mm, right=15mm, bottom=15mm,
	headsep=10mm, footskip=12mm}


\begin{document}
	\begin{tikzpicture}(10,3)
		\draw[ultra thick](2,0) --(10,0) -- (10,1.3) --(2,1.3) -- (2,0);
		\draw[fill=black](2,0)-- (0,.6) -- (2,.6) -- (2,0);
		\node at (6,.6) {\large solution sheet from \href{http://www.okuyakl.com}{www.okuyakl.com}};
	\end{tikzpicture}
	\vspace{0.5 cm}
	
	\noindent{\bf solution}\\
	\noindent{\bf Task 1.}\\
	$$|\overrightarrow{PQ}|=6 \quad \Rightarrow \quad \overrightarrow{PQ}=
	\left(\begin{array}{c} 6 \\ 0 \\ 0 \end{array} \right)\qquad \Rightarrow \quad \textnormal{e.g.} \quad P(1|1|1); \quad Q(7|1|1)$$
	
	\noindent{\bf Task 2.}\\
	Firstly, the vectors must be perpendicular to each other...
	$$\overrightarrow{AB} \circ \overrightarrow{AD}= \quad
	\left( \begin{array}{c} {1}\\ {2}\\ {2} \end{array} \right) \circ
	\left( \begin{array}{c} {2}\\ {1}\\ {-2} \end{array} \right) =2+2-4=0$$
	
	$$\overrightarrow{AB} \circ \overrightarrow{AE}= \quad
	\left( \begin{array}{c} {1}\\ {2}\\ {2} \end{array} \right) \circ
	\left( \begin{array}{c} {2}\\ {-2}\\ {1} \end{array} \right)=2-4+2=0 $$
	
	$$\overrightarrow{AD} \circ \overrightarrow{AE}= \quad
	\left( \begin{array}{c} {2}\\ {1}\\ {-2} \end{array} \right) \circ
	\left( \begin{array}{c} {2}\\ {-2}\\ {1} \end{array} \right)=4-2-2=0$$
	
	...and secondly, they have to be the same length:
	
	$$|\overrightarrow{AB}|=\sqrt{1+2^2+2^2} = 3 \qquad |\overrightarrow{AE}|=\sqrt{2^2+(-2)^2+1 } = 3 \qquad |\overrightarrow{AD}|=\sqrt{2^2+1+(-2)^2} = 3 ;$$
	
	\noindent{\bf Task 3. a)}\\
	\begin{minipage}{0.3\textwidth}
		\setlength{\unitlength}{0.5 cm}
		\begin{picture}(8,6)
			\put(1,1){\vector(1,0){5}}
			\put(1,1){\vector(1,2){1.5}}
			\put(2.5,4){\vector(1,0){5}}
			\put(6,1){\vector(1,2){1.5}}
			\put(0.5,0.5){O}
			\put(6.2,0.5){B}
			\put(7.8,4){C}
			\put(2,4){D}
		\end{picture}
	\end{minipage}
	\begin{minipage}{0.7\textwidth}
		The vector chain applies:
		$$\overrightarrow{OC}=\overrightarrow{OB}+\overrightarrow{BC}$$
		With $\overrightarrow{BC}=\overrightarrow{OD}$ it follows:
		$$\overrightarrow{OC}=\overrightarrow{OB}+\overrightarrow{OD}=
		\left( \begin{array}{c} {-3}\\ {4}\\ {0} \end{array} \right) +
		\left( \begin{array}{c} {-2}\\ {1}\\ {2} \end{array} \right) =
		\left( \begin{array}{c} {-5}\\ {5}\\ {2} \end{array} \right) \Rightarrow \quad
		C(-5|5|2)$$
	\end{minipage}
	\noindent{\bf Task 3. b)}\\
	The content of the base area $G$ is equal to the amount of the vector product:
	$$G=|\overrightarrow{OB} \times \overrightarrow{OD}|=
	\left| \left( \begin{array}{c} {-3}\\ {4}\\ {0} \end{array} \right) \times
	\left( \begin{array}{c} {-2}\\ {1}\\ {2} \end{array} \right)\right| =
	\left| \left( \begin{array}{c} {8}\\ {6}\\ {5} \end{array} \right)\right|=\sqrt{8^2+6^2+5^2 }=5\sqrt{5}$$
	\noindent{\bf Task 3. c)}\\
	\begin{minipage}{0.3\textwidth}
		\includegraphics[width=5 cm]{pyvec337}
	\end{minipage}
	\begin{minipage}{0.7\textwidth}
		The center of the base area $M_G$ has the coordinates: $$\overrightarrow{OM_G} = {1 \over 2} \cdot(\vec{0}+\overrightarrow{OC})={1 \over 2}\left( \begin{array}{c} {-5}\\ {5}\\ {2} \end{array} \right)=\left( \begin{array}{c} {-2,5}\\ {2,5}\\ {1} \end{array} \right)\Rightarrow \quad
		M_G(-2.5|2.5|1)$$
		The vector$\overrightarrow{M_GS}$, which connects the tip $S$ to the center, is perpendicular to the base and therefore runs in the direction of the vector
		$\overrightarrow{OB} \times \overrightarrow{OD}= \left( \begin{array}{c} {8}\\ {6}\\ {5} \end{array} \right)$. Since this already has the length $5 \sqrt{5}$, we can immediately adopt it as $$\overrightarrow{M_GS}$$.
	\end{minipage}
	The vector chain then applies to the location vector of the tip:
	$$\overrightarrow{OS}=\overrightarrow{OM_G}+\overrightarrow{M_GS}=\left( \begin{array}{c} {-2,5}\\ {2,5}\\ {1} \end{array} \right)+ \left( \begin{array}{c} {8}\\ {6}\\ {5} \end{array} \right)
		=\left( \begin{array}{c} {5,5}\\ {8,5}\\ {6} \end{array} \right)$$
		
		\noindent{\bf Task 3. d)}\\
		Point $N$ on the x-axis has the same $x_2$ coordinate as point $D$, so $N(0|1|0)$.
		The radius vector of the circle with its magnitude is:
		$$\overrightarrow{ND}=\left( \begin{array}{c} {-2}\\ {1}\\ {2} \end{array} \right)-\left( \begin{array} {c} {0}\\ {1}\\ {0} \end{array} \right)=
		\left( \begin{array}{c} {-2}\\ {0}\\ {2} \end{array} \right) \qquad |\overrightarrow{ND}|= \sqrt{4+4} =2\sqrt{2}$$
		
		\begin{minipage}{0.4\textwidth}
			\noindent{\bf Task 4. a)}\\
			Points $A$ and $D$ both lie in the $x_1x_3$ plane.\\
			\includegraphics[width=7 cm]{eder337}
		\end{minipage}
		\hspace{0.5 cm}
		\begin{minipage}{0.5\textwidth}
			\noindent{\bf Task 4. b)}\\
			A regular tetrahedron consists of four equilateral triangles. It has to be shown that:
			$$|\overrightarrow{AB}|=|\overrightarrow{AC}|=|\overrightarrow{AD}|$$
			$$|\overrightarrow{AB}|=\left|\left( \begin{array}{c} {-4}\\ {5}\\ {-3} \end{array} \right)\right| =
			\sqrt{16+25+9}=\underline{5\sqrt{2}}$$
			$$|\overrightarrow{AC}|=\left|\left( \begin{array}{c} {-3}\\ {5}\\ {4} \end{array} \right)\right|=
			\sqrt{9+25+16}=\underline{5\sqrt{2}}$$
			$$|\overrightarrow{AD}|=\left|\left( \begin{array}{c} {-7}\\ {0}\\ {1} \end{array} \right)\right|=
			\sqrt{49+1}=\underline{5\sqrt{2}}$$
			\end{minipage}\\
			
			It must also apply that the angle between the vectors mentioned is $60^\circ$:
			$$\angle(\overrightarrow{AB},\overrightarrow{AC})=\cos^{-1}{\left( \begin{array}{c} {-4}\\ {5}\\ {- 3} \end{array} \right) \circ \left( \begin{array}{c} {-3}\\ {5}\\ {4} \end{array} \right) \over \sqrt{ 50} \cdot \sqrt{50}}=\cos^{-1}{12+25-12 \over 50}=\cos^{-1}{1 \over 2} = 60^\circ$$
			$$\angle(\overrightarrow{AB},\overrightarrow{AD})=\cos^{-1}{\left( \begin{array}{c} {-4}\\ {5}\\ {- 3} \end{array} \right) \circ \left( \begin{array}{c} {-7}\\ {0}\\ {1} \end{array} \right) \over \sqrt{ 50} \cdot \sqrt{50}}=\cos^{-1}{28-3 \over 50}=\cos^{-1}{1 \over 2} = 60^\circ$$
			\noindent{\bf Task 4. c)}\\
			The center $M_T$ of the tetrahedron is:
			$$\overrightarrow{OM_T}={1 \over 4}\cdot(\overrightarrow{OA}+\overrightarrow{OB}+\overrightarrow{OC}+\overrightarrow{OD})={1 \over 4}\cdot \left[
			\left( \begin{array}{c} {4}\\ {0}\\ {3} \end{array} \right)+
			\left( \begin{array}{c} {0}\\ {5}\\ {0} \end{array} \right)+
			\left( \begin{array}{c} {1}\\ {5}\\ {7} \end{array} \right)+
			\left( \begin{array}{c} {-3}\\ {0}\\ {4} \end{array} \right)\right]=
			\left( \begin{array}{c} {0.5}\\ {2.5}\\ {3.5} \end{array} \right)
			\Rightarrow \quad M_T(0.5|2.5|3.5)$$
			We form the vectors $\overrightarrow{M_TA}$ and $\overrightarrow{M_TB}$:
			$$\overrightarrow{M_TA}=\left( \begin{array}{c} {3.5}\\ {-2.5}\\ {-0.5} \end{array} \right) \qquad
			\overrightarrow{M_TB}=\left( \begin{array}{c} {-0.5}\\ {2.5}\\ {-3.5} \end{array} \right) $$
			The angle between these two vectors is:
			$$\cos{\varphi}=\frac{ \left( \begin{array}{c} {3.5}\\ {-2.5}\\ {-0.5} \end{array} \right) \circ
			\left( \begin{array}{c} {-0.5}\\ {2.5}\\ {-3.5} \end{array} \right) }{\sqrt{3.5^2 +2.5^2+0.5^2}\cdot \sqrt{0.5^2+2.5^2+3.5^2}}=\frac{-6.25}{18.75}=-{1 \over 3} \Rightarrow
			\qquad \varphi \approx 109.5^\circ$$
			\noindent{\bf Task 4. d)}\\
			The surface area is four times the area of the triangle:
			$$O_T=4\cdot {1 \over 2} |\overrightarrow{AB} \times \overrightarrow{AD}|=2\cdot
			\left| \left( \begin{array}{c} {-4}\\ {5}\\ {-3} \end{array} \right) \times
			\left( \begin{array}{c} {-7}\\ {0}\\ {1} \end{array} \right)\right| = 2 \cdot
			\left| \left( \begin{array}{c} {5}\\ {25}\\ {35} \end{array} \right)\right|=2\cdot
			\sqrt{5^2+25^2+35^2}=50\sqrt{3} \approx 86.60\, {\rm FE}$$
			The volume is calculated using the spate product:
			$$V_T={1 \over 6}(\overrightarrow{AB}\times \overrightarrow{AD})\circ \overrightarrow{AC}={1 \over 6}
			\left( \begin{array}{c} {5}\\ {25}\\ {35} \end{array} \right)\circ \left( \begin{array}{c} {-3}\ \ {5}\\ {4} \end{array} \right)={250 \over 6} \approx 41.7 \, {\rm VE} $$
			\newpage
			\noindent{\bf Task 4. e)}\\
			\begin{minipage}{0.7\textwidth}
			The triangle $ABC$ is formed by centric stretching with the center $A$ and the stretching factor
			$k={1 \over 2}$ mapped to the triangle $AM_1M_2$. All dimensions of the small triangle are then half the size of the original triangle, the area is then
			$$A_d=\left({1 \over 2}\right)^2 \cdot A_D={1 \over 4} \cdot A_D= 5.41 \, FE$$
			The volume of the small tetrahedron is
			$$V_t=\left({1 \over 2}\right)^3\cdot V_T={1\over 8}\cdot V_T ={125\over 24}\approx 5.21 \, {\rm VE} $$
			\end{minipage}
			\begin{minipage}{0.3\textwidth}
			\setlength{\unitlength}{0.8 cm}
			\begin{picture}(7,7)
				\put(1,3){\line(2,-1){5}}
				\put(1,3){\line(2,1){5}}
				\put(1,3){\line(2,-1){5}}
				\put(3.5,4.25){\line(0,-1){2.5}}
				\put(6,5.5){\line(0,-1){5}}
				\put(0.3,3.5){$A$}
				\put(3.1,1.1){$M_1$}
				\put(3.1,4.7){$M_2$}
				\put(6.3,0.5){$B$}
				\put(6.3,5.6){$C$}
			\end{picture}
			\end{minipage}
			
			\noindent{\bf Task 4. f)}\\
			The volume of the truncated pyramid is:
			$$V_{S}=V_T-V_t= {250 \over 6} \cdot \left( 1 - {1\over 8} \right) \approx 36.46 \,{\rm VE} $$
			The surface area of the frustum is equal to that of the large tetrahedron minus three small triangles plus one small triangle as the cutting surface:
			$$O_S=O_T-3\cdot A_d + A_d = O_T - 2 A_d= 50\sqrt{3} -2\cdot 5.41 = 75.78 \, {\rm FE}$$
			
			\noindent{\bf Task 5.}\\
			\begin{minipage}{0.3\textwidth}
			\setlength{\unitlength}{0.5 cm}
			\begin{picture}(10,8)
				\put(1,1){\includegraphics[width=4 cm]{parl337}}
				\put(4,1){$\vec{a}$}
				\put(1.5,5){$\vec{b}$}
				\put(3,2.5){$\vec{a}+\vec{b}$}
				\put(3.5,6.3){$\vec{a}-\vec{b}$}
			\end{picture}
			\end{minipage}
			\begin{minipage}{0.7\textwidth}
			A parallelogram is considered, which is spanned by the vectors $\vec{a}$ and $\vec{b}$. Then the diagonals are $\vec{a}+\vec{b}$ and $\vec{a}-\vec{b}$. The squares over the vectors have the areascontent $|\vec{v}|^2 = \vec{v}\circ \vec{v}$. The following applies to the area $A_Q$ of the diagonal squares:
			$$A_Q=(\vec{a}+\vec{b})\circ (\vec{a}+\vec{b}) + (\vec{a}-\vec{b})\circ (\ vec{a}-\vec{b})=
			\vec{a}^2 +2\cdot \vec{a}\circ \vec{b} + \vec{b^2} + \vec{a}^2 -2\cdot \vec{a}\circ \vec{b} +
			\vec{b}^2= 2\vec{a}^2+2\vec{b}^2$$
			This corresponds exactly to the area of the squares over the two $\vec{a}$ and the two $\vec{b}$ vectors.
			\end{minipage}
			
			
			\noindent{\bf Task 6. a)}\\
			The coordinates of the points in the $x_1x_2$ plane are:
			$$A(2|-3|0) \qquad B(1|4|0) \qquad C(-3|-4|0)$$
			\noindent{\bf Task 6. b)}\\
			The area is provided by the vector product. Original triangle:
			$$A_D={1\over 2}\cdot |\overrightarrow{AB} \times \overrightarrow{AC}|={1\over 2}\cdot
			\left| \left( \begin{array}{c} {-1}\\ {7}\\ {2} \end{array} \right) \times
			\left( \begin{array}{c} {-5}\\ {1}\\ {-2} \end{array} \right)\right|={1\over 2}
			\left| \left( \begin{array}{c} {-16}\\ {-12}\\ {34} \end{array} \right) \right|
			={1\over 2}\sqrt{16^2+12^2+34^2}\approx 19.72 \, {\rm FE}$$
			Projected triangle:
			$$A_D'={1\over 2}\cdot |\overrightarrow{A'B'} \times \overrightarrow{A'C'}|={1\over 2}\cdot
			\left| \left( \begin{array}{c} {-1}\\ {7}\\ {0} \end{array} \right) \times
			\left( \begin{array}{c} {-5}\\ {1}\\ {0} \end{array} \right)\right|={1\over 2}
			\left| \left( \begin{array}{c} {0}\\ {0}\\ {34} \end{array} \right) \right|
			= {1\over 2} \cdot 34 =17\, {\rm FE}$$
			The area of the projected triangle is around
			$$(1-{17 \over 19.72})\cdot 100\% = 13.8\%$$
			smaller than that of the original.
			
			\noindent {\bf Task 7. a)}\\
			The point $A$ on the sphere satisfies the vector equation:
			$$
			\renewcommand{\arraystretch}{1.5}
			\begin{array}{rcll}
			[\overrightarrow A -\overrightarrow M]^2 &=& r^2 &(I)\\
			\\
			\left[ \left( \begin{array}{c} 8 \\ 2 \\ 5 \end{array}\right)-
			\left( \begin{array}{c} 2 \\ -1 \\ 7 \end{array}\right)\right]^2&=& \\
			6^2 + 3^2 + (-2)^2 &=& 49
			\end{array}
			$$
			So the sphere has radius $r=\sqrt{49}=7$. We now put the point $B_k$ in
			(I) a:
			$$
			\renewcommand{\arraystretch}{1.5}
			\begin{array}{rcll}
			[\overrightarrow B_k -\overrightarrow M]^2 &=& 49 \\
			\\
			\left[ \left( \begin{array}{c} 5 \\ 5 \\ k \end{array}\right)-
			\left( \begin{array}{c} 2 \\ -1 \\ 7 \end{array}\right)\right]^2&=&\\
			3^2 + 6^2 +(k-7)^2 &=& 49 &|-45 \\
			(k-7)^2 &=& 4 &|\sqrt{\dots}\\
			k-7 &=& \pm 2 \\
			\textnormal{\underbar{Case 1:}}\qquad k &=& 9 \Rightarrow & B_9(5|5|9)\\
			\textnormal{\underbar{Case 2:}}\qquad k &=& 5 \Rightarrow & B_5(5|5|5)
			\end{array}
			$$
			
			\noindent {\bf Task 7. b)}\\
			We form the vectors $\overrightarrow {B_9M}$ and $\overrightarrow{AM}$:
			$$\overrightarrow{AM}= \left( \begin{array}{c} 6 \\ 3 \\ -2 \end{array}\right) \qquad
			\overrightarrow{B_9M}= \left( \begin{array}{c} 3 \\ 6 \\ 2 \end{array}\right)$$
			The area of a triangle spanned by two vectors is equal to half the value of the vector product:
			$${1\over 2}\left| \left( \begin{array}{c} 6 \\ 3 \\ -2 \end{array}\right) \times
			\left( \begin{array}{c} 3 \\ 6 \\ 2 \end{array}\right)\right|=
			{1\over 2}\left| \left( \begin{array}{c} 18 \\ -18 \\ 27 \end{array}\right)\right| =
			{1\over 2}\sqrt{18^2+18^2+27^2} \approx 18.6 \, {\rm FE}$$
			
			\begin{center}
			\includegraphics[width=7 cm]{../../viecher/eendcomic.pdf}
			
			Here you can go back to the \href{https://www.okuyakl.de/math/m11vecL337/ae337.pdf}{task sheet}
			\end{center}
			
		\end{document}